
\RequirePackage[T1]{fontenc}
\documentclass[conference,a4paper]{IEEEtran}
\usepackage{amsmath,amssymb,booktabs,graphicx,cite,url}
\usepackage[hidelinks]{hyperref}
\graphicspath{{figures/revision/}}
\newcommand{\MainClosed}{6.70}
\newcommand{\MainDegraded}{97.53}
\newcommand{\MainQualified}{6.61}
\newcommand{\MainForced}{90.08}
\newcommand{\MainForcedLo}{89.20}
\newcommand{\MainForcedHi}{90.97}
\newcommand{\MainPayloadLoss}{92.24}
\newcommand{\MainNoAttack}{97.53}
\newcommand{\MainQualifiedDelta}{-0.09}
\newcommand{\MainQualifiedDeltaLo}{-0.11}
\newcommand{\MainQualifiedDeltaHi}{-0.07}

\newcommand{\IEEEauthorblockN}[1]{#1\\}
\newcommand{\IEEEauthorblockA}[1]{#1}
\newcommand{\IEEEauthorrefmark}[1]{\textsuperscript{#1}}
\begin{document}
\title{Fail-Closed or Fail-Degraded? A Simulation Study of Forced Assurance Loss in Contested Multi-Hop QKD Networks}
\author{
\IEEEauthorblockN{Lakshya Jain\IEEEauthorrefmark{1}, Siddhi Arora\IEEEauthorrefmark{2}, Divya Mishra\IEEEauthorrefmark{3}}
\IEEEauthorblockA{\IEEEauthorrefmark{1}\textit{Maharaja Agrasen Institute of Technology}, Delhi, India, 10jainlakshya@gmail.com\\
\IEEEauthorrefmark{2}\textit{Noida Institute of Engineering and Technology}, Greater Noida, India, Sid.arora1701@gmail.com\\
\IEEEauthorrefmark{3}\textit{G.L. Bajaj Institute of Technology and Management}, Greater Noida, India, divya\_rbl2@yahoo.com}}
\maketitle

\begin{abstract}
Stopping quantum key generation need not immediately exhaust a trusted-relay network's stored keys. We study how finite reserves, interruption duration and fallback policy jointly determine timely key service and assurance loss. The evaluation separates generation failure, classical transport failure and online key-establishment failure, and compares policies on identical request and environmental realizations. In 360 simulations over five network structures, a 24-hour generation cut with only 2\% of usable capacity initially filled preserves high service under computational fallback while sharply reducing floor-compliant service. In the eight-node formation, fail-closed service is \MainClosed\%, compared with \MainDegraded\% under fallback; compliant fallback service is only \MainQualified\%. The fraction of offered, ITS-key-entitled requests served below that entitlement is \MainForced\% (95\% interval \MainForcedLo--\MainForcedHi\%). We also evaluate an interpretable policy learned under deterministic class-specific exception budgets, using separate training, validation and test realizations. The results characterize a modeled key-supply trade-off, rather than payload throughput or independently validated field security.
\end{abstract}
\begin{quote}Index Terms---
Quantum key distribution, trusted relays, key management, post-quantum cryptography, assurance metrics, constrained policy search.
\end{quote}

\section{Introduction}
When a quantum key distribution (QKD) network cannot supply a requested key, a controller can wait or refuse the request, or use a computational alternative. ITU-T Y.3803 permits policy-controlled symmetric substitution during key relay and records the method used~[1]. Tysowski \emph{et al.} describe both fallback and blocking~[2]; Spooren \emph{et al.} implement a layered fail-safe design with finite refresh and failure timers~[3]. Both choices are therefore represented in standards and system designs. Their evaluation depends on what is counted as service, which failures affect existing keys, and how much usable stock remains.

This matters for a relay formation carrying command-key requests, an expensive one-time-pad class and less restrictive logistics traffic. A successful key request does not imply that its original assurance requirement was met. Conversely, stopping a generator does not imply that previously accepted, synchronized and uncompromised keys are unusable. An instantaneous generation-cut model can therefore confuse lost replenishment with lost service. A reservoir-retaining model needs an explicit alternative for implementations that intentionally purge working keys after a failure.

Our contribution is a reproducible empirical evaluation of this joint decision: finite reserves and interruption duration, heterogeneous link-suppression budgets, entitlement-conditioned loss, and learned allocation under explicit degradation permissions. We report timely and compliant key service with the same arrival denominator, retain a separate ITS-key entitlement metric, and include payload and declared class-importance weightings. A cumulative cut-budget inequality provides a necessary resource check; small exact cases demonstrate why it is not generally an achievable service count. The learning study tests whether a compact controller adds value beyond a greedy policy under the same runtime constraints. Neither the inequality nor the use of finite buffers, fallback or safety filtering is claimed as a new general theory.

\section{Related Work}
Adaptive graded key delivery already appears in QuLore~[4], and tabular Q-learning has been applied to resource-aware trusted-relay routing~[5]. Our learned action is a requested assurance grade, rather than a next-hop decision. Routing is held within a small, explicitly defined path-selection implementation; we do not claim to reproduce or outperform a published routing algorithm.

Three recent studies further narrow the contribution. Opportunistic QKD analyzes a single-link reservoir, temporally correlated supply and a depletion/recovery horizon~[6]. QuIKS combines end-to-end buffer sizing with reactive control and device-trace evaluation~[7]. Zhu studies finite buffers, preconfigured fallback fractions and class-sensitive security output across several network scenarios~[8]. These works establish that depletion time, buffering and weighted service value are prior art. Our evaluation instead combines targeted generation interruptions, declared per-link cost, paired excess entitlement loss and a learned class-exception allocator. This is a focused comparison of the inspected methods, not an exhaustive absence claim.

Runtime action filtering also has established foundations~[9]. Here the monitor enforces a specific prefix-count permission contract. Its guarantee follows from integer accounting, not from the learning procedure, and does not imply composable cryptographic security.

\section{Model and Assurance Accounting}\label{sec:model}
\subsection{Stored key and failure surfaces}
A trusted-relay graph has one shared key store per undirected link. Relaying a session key debits every traversed link once. For link $e$, stored stock obeys


\[
B_e(t)=B_e(0)+G_e(t)-C_e(t)-O_e(t)-P_e(t),
\]

WORD-EQUATION-NUMBER:1


where $G_e$ is gross incoming distilled-key supply before buffer overflow, $C_e$ useful relay consumption, $O_e$ overflow discard, and $P_e$ intentional working-key purge. Usable stock is $U_e=(B_e-b_{\min})_+$. Admission checks every hop before any debit; partial relay transactions and in-flight failures are outside this abstraction.

The retained-key case assumes stored keys remain synchronized, unexpired and uncompromised when generation stops. The alternative discards spendable working stock at a transition to failed generation, while retaining the authentication reserve. It is an explicitly conservative policy, not a claim that every implementation behaves identically. Standards distinguish key generation and stored-key management~[1]; documented implementations can purge keys under particular shutdown or synchronization failures~[2].

Generation suppression leaves classical relay transport available. Transport failure prevents key relay and modeled session service even if key material is pre-positioned. An isolated online key-establishment failure is a third surface: it can prevent a KEM-based fallback without blocking a valid pre-positioned symmetric key. We model these surfaces separately. Payload transmission bandwidth, cipher execution and an actual KEM protocol are not simulated.

Nominal generation is a parameterized input. For the QBER sensitivity, with $e_0=0.03$ and reconciliation factor $f=1.16$, we use the reduced response


\[
\frac{R_e(e)}{R_{e,0}}=\left[\frac{\max\{0,1-(1+f)H_2(e)\}}{1-(1+f)H_2(e_0)}\right]^{\gamma}
\]

WORD-EQUATION-NUMBER:2


below $e_{\mathrm{abort}}$, and zero above it; the ratio is capped at one. The engineering threshold is 0.075, with 0.11 as a sensitivity. The entropy numerator can become zero before a selected abort threshold, so the operational threshold alone is not a security-rate guarantee. The exponents 1 and 3.4 are illustrative response shapes, not a cross-modality calibration or a finite-key security proof.

\subsection{Grades, requests and debits}
The four grades are policy labels: G1 reserves one-time-pad key equal to payload size; G2 reserves a 256-bit QKD-delivered symmetric key; G3 reserves 128 QKD bits as a modeled hybrid-key ingredient; G4 uses a computational fallback with no quantum-key debit. They are an ordered preference in this experiment, not a universal ordering of cryptographic strength. In particular, G2 denotes ITS \emph{key delivery}, not information-theoretically secure payload encryption. No hybrid-combiner security claim is inferred from its accounting label.

For a nonzero quantum-key demand $D_g$, the per-hop debit is


\[
q_g=D_g+240+127n_{\mathrm{auth}}\quad\text{bits}.
\]

WORD-EQUATION-NUMBER:3


The metadata allocation is an explicit choice within the key-file abstraction of Y.3803~[1]. The tag size is motivated by Walenta \emph{et al.}~[10]; reusable-hash authentication charges are discussed by Fung \emph{et al.}~[11]. These sources do not establish a universal number of authenticated messages per relayed request. We use $n_{\mathrm{auth}}=7$ as a conservative relay-overhead setting relative to the optimistic value 2, and also evaluate 40. The corresponding G2 debits are 1385, 750 and 5576 bits per hop. This is a simplified cost model, not protocol-level calibration.

C2, ISR and LOG requests have respective payloads 2048, 144000 and 320000 bits; floors G2, G1 and G3; deadlines 5, 30 and 120 seconds; and strict scheduling priorities in that order. Independent Poisson arrival rates are 0.05, 0.002 and 0.03 per second in the reduced-load case. These are declared workload scenarios. The required mean key debit is 397.218 bits/s across a cut traversed once per request, compared with 480 bits/s for two nominal 240-bit/s branches. This is a nominal capacity check, not a queue-stability proof. ISR arrival-rate factors 0.03, 0.1 and 1 relative to the historical 0.02/s workload separate lighter load from overload.

\subsection{Common-denominator outcomes}\label{sec:metrics}
For all $N$ requests arriving in a fixed observation window, let $S$ receive key service within their deadlines and $D\subseteq S$ receive a grade below their class floor. Define


\[
A=\frac{|S|}{N},\qquad Q=\frac{|S\setminus D|}{N},\qquad L=\frac{|D|}{N}.
\]

WORD-EQUATION-NUMBER:4


Thus $A=Q+L$; blocking is $1-A$. This identity makes fallback's service advantage distinguishable from floor compliance.

Let $E$ be requests entitled to G1 or G2, and $D_{\mathrm{ITS}}$ the members served at G3 or G4. Offered-request forced assurance loss is $F=|D_{\mathrm{ITS}}|/|E|$. ITS-delivered, below-ITS-delivered and blocked fractions sum to one on this common entitled-arrival cohort. The conditional served-request fraction $F_{\mid S}$ is retained in the released tables, but is not substituted for $F$. LOG is excluded from the ITS-entitled denominator rather than counted as an attack success by default.

Payload exposure uses the corresponding sum of associated payload bits instead of counts. It measures a declared exposure proxy, not observed compromise or calibrated harm. Class-importance vectors $(1,1,1)$, $(5,2,1)$ and $(10,3,1)$ provide an additional transparent reweighting. Attack-attributable change is estimated as a paired attacked-minus-unattacked difference under the same exogenous environment; no zero baseline is assumed.

\subsection{A finite cut-budget check}\label{sec:bound}
For any cut separating a request's endpoints, every compliant relay must consume key on at least one crossing edge. Let $\mathcal C_T$ contain such compliant requests completed by time $T$, and $q_s$ be their minimum per-hop debit. Then


\[
\sum_{s\in\mathcal C_T}q_s
\leq\sum_{e\in\delta(U)}\left[U_e(0)+G_e(T)\right].
\]

WORD-EQUATION-NUMBER:5


Discard, competing traffic and repeated cut crossings only tighten this upper bound. It follows by summing conservation over crossing stores. A generation cut removes replenishment, not the initial term.

The pooled bound need not be attainable. With three disjoint branches, 1900 spendable bits per edge and a 750-bit debit, it permits seven requests, but the simulator serves only six: each branch retains an unusable 400-bit fragment. Four 749-bit stores permit three requests under the pooled relaxation and none in the actual graph. Exact divisible-stock and purge cases also pass. These fixtures check implementation and the bound's scope; they do not validate hardware behavior.

\section{Policies and Experimental Protocol}\label{sec:protocol}
\subsection{Compared control policies}
Fail closed waits for the authorized floor and blocks at deadline. Unrestricted fallback attempts progressively weaker grades when that floor cannot be met. Our shortest-path baseline chooses the first transport-usable minimum-hop candidate. The key-aware implementation minimizes bottleneck estimated refill time plus a 0.5-second penalty per hop. Candidate enumeration is limited to four simple paths within two extra hops of the shortest. The comparison concerns these implementations.

The learned controller uses 18 states: three classes, three floor-feasibility categories, and two deadline-urgency categories. Feasibility distinguishes affordable now, refill estimated within the remaining deadline, and neither; urgency means at most 20\% of the deadline remains. Estimated refill uses current rates and ignores future competing arrivals. The five proposed actions are the class floor, strongest affordable fallback, defer, G4 and G2.

All learned and budget-greedy actions pass through the same monitor. If $N_c(t)$ class-$c$ requests have arrived and $V_c(t)$ below-floor services have been committed, a proposed exception is permitted only when


\[
V_c(t)+1\leq\lfloor\epsilon_cN_c(t)\rfloor,
\]

WORD-EQUATION-NUMBER:6


with $(\epsilon_{\mathrm{C2}},\epsilon_{\mathrm{ISR}},\epsilon_{\mathrm{LOG}})=(0,0.1,0.1)$. Failed attempts spend no allowance. Since arrivals only increase the right-hand side and every committed exception is checked before use, induction gives the prefix bound. Completed exceptions are a subset of committed exceptions. This is an arrival-count permission guarantee within the simulator, not a learned cryptographic guarantee. Strict-floor and greedy policies are explicit baselines.

Categorical cross-entropy policy search~[12] uses ten unique tables per generation, four generations and two optimizer restarts. Every generation includes the strict, greedy and incumbent tables. Short-nominal training uses six-hour rollouts; full-nominal uses 24 hours; full-robust uses a mean over nominal, generation-cut and correlated low-rate conditions. Each candidate receives six rollouts. Validation selects the restart on twelve separate condition/seed cells. Final testing uses twelve previously unused seeds over seven fixed conditions, including unseen topologies. The raw tables, complete search histories, objective $J=A-0.75L$, validation rule and test plan are frozen before final testing. A selected table is the best found under this budget, not an optimum.

\subsection{Scenarios, uncertainty and verification}
The graph families are an eight-node, nine-link formation; an eight-node ring; a $3\times3$ grid; and three or four parallel three-hop branches. The latter have $(8,9)$ and $(10,12)$ nodes/edges. Individual formation terminals are leaves, with cut size one; isolating all formation terminals requires two edges. The other source-terminal cut sizes are two, two, three and four. These different notions of cut are not conflated.

Formation rates are 200--400 bits/s on eight free-space links and 300 kbits/s on one fiber; the other graphs use 240-bit/s links. Capacities, threshold and reserve are selected legacy QKDNetSim example parameters~[13]: 52428800, 11324620 and 1048576 bits. Capacity is 50 Mibit, not a measured hardware specification. Our overflow model clamps continuous stock and discards excess; it is not byte-for-byte equivalence to a discrete block buffer. Initial spendable fractions are 0.02, 0.1 and 1; the main stressed initial stock is 1027604.48 usable bits per edge.

Link-suppression cost is one for free-space links and ten for fiber, with ratios one and one hundred as sensitivities. These are relative scenario weights, not measured adversary expenditure. A weighted source-to-all-terminals cut is a structural reference, not an optimal finite-horizon attack. Independent generation and transport failures use declared mean up/down times of 3600/120 seconds. The correlated-generation sensitivity gives all free-space generators a common 900/600-second up/down process. Exogenous draws are independent of policies and attack choices.

Simulation uses one-second steps and a 24-hour arrival window without a steady-state claim. Pending requests are followed to their deadlines; an interruption scheduled to end at 24 hours can therefore release before the last completion. Equation~5 uses the actual completion horizon when applied to such a cohort. Twelve seeds support the main grid, eight the sensitivity and fallback comparisons, and six the threshold sweeps. Intervals are marginal 95\% percentile bootstrap intervals over independent seeds, with 20000 resamples. They represent simulation uncertainty conditional on the parameters and, for learning, the selected table.

Verification comprises conservation, reserve/deadline invariants, deterministic resource fixtures, trace-matched counterfactuals and independent prefix-budget audits. Arithmetic drain checks and broad published capacity ranges do not constitute independent network validation. Aggregate measurements from unlike relay and application protocols do not calibrate every accounting term. No hardware, field-test or composable-security validation is claimed.

\section{Results}\label{sec:results}
\subsection{Availability is not compliant service}
At the main low-reserve setting, suppressing the weighted generation cut leaves formation fail-closed service at \MainClosed\%, compared with \MainNoAttack\% without attack. Fallback serves \MainDegraded\% of arrivals, but only \MainQualified\% meet their original floor. Its compliant-service difference from fail closed is \MainQualifiedDelta\ percentage points, with interval $[\MainQualifiedDeltaLo,\MainQualifiedDeltaHi]$. The large availability difference therefore cannot be interpreted as preserved authorized protection.

\begin{figure}[t]
\centering
\includegraphics[width=3.486924in]{C:/Users/10jai/OneDrive/Desktop/ResearchPaper/build/revised_word_preview/staging/assets/fig1_comparable_service.png}
\caption{Fig.~1. Timely and floor-compliant key service under full generation-cut suppression, 24-hour arrival cohorts and 2\% initial usable stock. Bars are 95\% seed-bootstrap intervals ($n=12$). The horizontal scales differ.}

\end{figure}

\begin{table}[t]
\caption{TABLE~I\\ Formation class outcomes under the two-edge generation cut. Entries are percent of class arrivals; 95\% intervals appear below the means, from 12 seeds.}

\centering\footnotesize
\setlength{\tabcolsep}{3pt}
\begin{tabular}{@{}lccc@{}}
\toprule
Class & Closed $Q$ & Fallback $A$ & Fallback $Q$\\
\midrule
C2 & {6.67\newline (5.70, 7.61)} & {96.72\newline (96.52, 96.89)} & {6.65\newline (5.68, 7.59)}\\[3pt]
ISR & {5.67\newline (5.41, 5.97)} & {97.50\newline (96.66, 98.32)} & {5.67\newline (5.41, 5.97)}\\[3pt]
LOG & {6.83\newline (5.91, 7.67)} & {98.88\newline (98.74, 99.01)} & {6.60\newline (5.73, 7.41)}\\[3pt]
\bottomrule
\end{tabular}

\end{table}

Figure~1 shows the same distinction across all five graphs. Table~I prevents an aggregate score from hiding a class's admission failure. Formation forced loss is \MainForced\% of offered ITS-entitled requests, whereas below-floor payload exposure is \MainPayloadLoss\% of all associated payload bits. The different values reflect different denominators and requirements, not conflicting estimates.



\section{Discussion and Limitations}\label{sec:limits}
The model separates key-supply continuity from compliance with the requested key grade. It assumes trusted relays, accepted-key integrity, synchronized stores, instant state information and atomic debits. Compromised relays or detector attacks that invalidate key material are outside its scope. Generation and traffic processes are declared experimental conditions; they do not reconstruct a measured tactical deployment. The candidate-path limit, fixed class payloads, controller observability and simplified computational fallback all constrain external validity. Payload exposure and class weights cannot identify the harm of an actual compromise.

Finite buffers, depletion analysis and runtime shielding have established precedents. The practical question here is their interaction with explicit entitlement, interruption cost and a constrained learned controller. A favorable scalar utility does not waive a class requirement, and a low forced-loss value can result from blocking rather than preserving service. Both outcomes must remain visible. Training-horizon comparisons concern the registered protocols, including their different simulated exposure and state coverage; they do not isolate depletion as the unique cause of a learned-policy difference.

\section{Conclusion}
For finite-reservoir trusted-relay QKD, a generation cut, a key-store policy and a service policy are separate modeling decisions. Across the tested graphs, fallback can preserve timely key service while providing little additional floor-compliant service. Common arrival denominators and explicit class outcomes make that distinction measurable. The accompanying cut-budget fixtures, parameter sweeps and constrained policy study provide a reproducible basis for evaluating the trade-off under stated assumptions.

\section*{Reproducibility and Acknowledgement}
The artifact includes fixed configurations, raw per-seed outcomes, source hashes, numerical-table and figure generators, tests, policy-search histories and a validation-selected policy freeze. Historical outputs are preserved separately. The authors used a large language model to assist coding, analysis and drafting; the reported numerical evidence is generated by the accompanying executable experiments.

\section*{References}

[1] ITU-T, \emph{Quantum key distribution networks---Key management}, Recommendation Y.3803, Dec. 2020.

[2] P. K. Tysowski, X. Ling, N. L\"utkenhaus, and M. Mosca, ``The engineering of a scalable multi-site communications system utilizing quantum key distribution (QKD),'' \emph{Quantum Sci. Technol.}, vol. 3, 024001, 2018, doi: 10.1088/2058-9565/aa9a5d.

[3] P. Spooren, A. Neuhold, S. Ramacher, and T. Huehn, ``PQC-enhanced QKD networks: A layered approach,'' in \emph{Proc. IEEE QCNC}, 2026, doi: 10.1109/QCNC69040.2026.00060.

[4] A. Sanz \emph{et al.}, ``QuLore: An adaptive security framework to extend quantum-safe communications to real-world networks,'' arXiv:2511.22416v2, 2026.

[5] Y. Hao \emph{et al.}, ``Q-learning for resource-aware and adaptive routing in trusted-relay QKD network,'' \emph{Photonics}, vol. 12, no. 10, 969, 2025, doi: 10.3390/photonics12100969.

[6] S. Chaudhary, P. Munar Vallespir, A. Viladomat Jasso, and J. Noetzel, ``Opportunistic QKD: Exploiting idle capacity of classical WDM systems,'' arXiv:2604.12982v2, 2026.

[7] Y. Chen \emph{et al.}, ``QuIKS: Near-zero latency key supply with adaptive buffering for resource-efficient quantum key distribution networks,'' arXiv:2604.09144v1, 2026.

[8] Z. Zhu, ``Techno-economic feasibility analysis of quantum key distribution for power-system communications,'' arXiv:2510.15248v1, 2025.

[9] M. Alshiekh \emph{et al.}, ``Safe reinforcement learning via shielding,'' arXiv:1708.08611v2, 2017.

[10] N. Walenta \emph{et al.}, ``A fast and versatile QKD system with hardware key distillation and wavelength multiplexing,'' \emph{New J. Phys.}, vol. 16, 013047, 2014, doi: 10.1088/1367-2630/16/1/013047.

[11] C.-H. F. Fung, X. Ma, and H. F. Chau, ``Practical issues in quantum-key-distribution post-processing,'' \emph{Phys. Rev. A}, vol. 81, 012318, 2010, doi: 10.1103/PhysRevA.81.012318.

[12] R. Rubinstein, ``The cross-entropy method for combinatorial and continuous optimization,'' \emph{Methodol. Comput. Appl. Probab.}, vol. 1, pp. 127--190, 1999, doi: 10.1023/A:1010091220143.

[13] QKDNetSim, ``SECOQC example and key-buffer implementation,'' \emph{qkdnetsim-v1-old}, commit 394bfed, 2024. [Online]. Available: \url{https://github.com/QKDNetSim/qkdnetsim-v1-old}

\end{document}
